% s-160-p160-m002; original PDF 160
\Aphorism{s-160-p160-m002}{245}
Sich an das atomistische Glück des Tags, der Stunde, der Minute verlieren können, setzt voraus, dass man nicht mit Hoffnungen und Erinnerungen über Jahre hinaus belastet sei: man würde sonst zu schwer, zu unlenkbar, zu „gerichtet“, zu reichlich mit Bewegungsenergie begabt — man würde durch die Erlebnisse gar zu mikroskopischer Gestalt ohne Aufenthalt hindurchfliegen wie das Geschoss durch einen Mückenschwarm. \Foreign{Carpe diem!} Wenn das Sieb auf Lebensglück gestellt ist, fällt Tag- und Stundenglück ungenossen heraus. Was liegt daran, ob ich an blühenden Auen oder staubgebräunten Trümmern vorüberwandere, wenn ich doch vorüber muss? Desshalb verbieten sich Weise, ein grosses Ziel und einen langen Weg zu haben, um eben nirgends vorüber zu müssen, um wie die Kinder den Augenblick ohne Vor- und Rückschau zu geniessen und überall das kleine Glück des Tages zu pflücken.
\par

